\documentclass[
    language=english,
    doctype=research-proposal,
    institution=none,
    titlestyle=book,
]{../../omnilatex}

\RequirePackage{config/document-settings}
\addbibresource{../../bib/bibliography.bib}

\begin{document}

\maketitle

\frontmatter
\chapter{Abstract}
This proposal outlines a three-year programme to develop scalable quantum error
correction (QEC) codes for near-term NISQ devices.  By leveraging topological
stabiliser codes and machine-learning--assisted decoding, we aim to reduce the
logical error rate by two orders of magnitude while keeping overhead within
feasible limits.

\tableofcontents

\mainmatter

\chapter{Introduction}
Quantum computing promises exponential speed-ups in chemistry, optimisation,
and cryptography~\cite{einstein1905}.  However, current NISQ devices are limited
by decoherence and gate errors.  Quantum error
correction~\cite{diracPrinciplesQuantumMechanics1981} is the path to fault
tolerance, but existing schemes impose prohibitive overheads.  This proposal
targets a practical QEC architecture bridging theory and near-term hardware.

\chapter{Background and Literature Review}
QEC has its foundations in the work of Shor and Steane, who demonstrated that
quantum information can be protected via entangled encoding.  The surface
code~\cite{knuthFundamentalAlgorithms1973} is a leading candidate due to its
local stabiliser measurements and high error threshold ($\approx 1\%$).
Bird et al.~\cite{birdNaturalLanguageProcessing2009} highlight the growing role
of data-driven decoding.  Two challenges remain: reducing physical qubits per
logical qubit, and achieving real-time decoding at microsecond timescales.

\chapter{Research Objectives}
The goal is a QEC system with logical error rate below $10^{-6}$ using fewer
than 1\,000 physical qubits per logical qubit:
\begin{enumerate}
    \item Design tailored topological codes for realistic noise models.
    \item Develop a neural-network decoder meeting real-time latency constraints.
    \item Validate through large-scale simulation and a 72-qubit demonstration.
\end{enumerate}

\chapter{Methodology}

\section{Code Design}
We construct $[[n,k,d]]$ stabiliser codes parameterised by hardware noise bias.
The code distance must satisfy
\begin{equation}
    d \geq \left\lceil \frac{\log(1/\epsilon_{\mathrm{target}})}
    {\log(1/p_{\mathrm{phys}})} \right\rceil
    \label{eq:distance}
\end{equation}
where $\epsilon_{\mathrm{target}}$ is the target logical error rate.
Table~\ref{tab:codes} summarises preliminary parameters.

\begin{table}[htbp]
    \centering
    \caption{Preliminary code parameters.}
    \label{tab:codes}
    \begin{tabular}{@{}rrrr@{}}
        \toprule
        $p_{\mathrm{phys}}$ & $n$ & $k$ & $d$ \\
        \midrule
        $10^{-3}$ & 489  & 1 & 11 \\
        $10^{-3}$ & 1089 & 2 & 15 \\
        $10^{-2}$ & 1449 & 1 & 19 \\
        \bottomrule
    \end{tabular}
\end{table}

\section{Neural-Network Decoder}
We train a CNN to map syndrome measurements to correction
operations~\cite{birdNaturalLanguageProcessing2009}.  Figure~\ref{fig:decoder}
illustrates the pipeline.

\begin{figure}[htbp]
    \centering
    \begin{tikzpicture}[
        box/.style={draw, thick, rounded corners, minimum height=1cm,
                     minimum width=2cm, fill=blue!10, font=\small},
        arr/.style={->, thick, >=stealth},
        node distance=1.2cm,
    ]
    \node[box] (input) {Syndrome\\Input};
    \node[box, right=of input] (conv1) {Conv2D\\$3{\times}3$};
    \node[box, right=of conv1] (conv2) {Conv2D\\$3{\times}3$};
    \node[box, right=of conv2] (pool) {MaxPool};
    \node[box, right=of pool] (fc) {Dense\\Layer};
    \node[box, right=of fc] (output) {Correction\\Output};
    \draw[arr] (input) -- (conv1);
    \draw[arr] (conv1) -- (conv2);
    \draw[arr] (conv2) -- (pool);
    \draw[arr] (pool) -- (fc);
    \draw[arr] (fc) -- (output);
\end{tikzpicture}
    \caption{Decoder pipeline from syndrome input to correction output.}
    \label{fig:decoder}
\end{figure}

\chapter{Timeline and Work Plan}
Table~\ref{tab:gantt} shows the projected 36-month timeline.

\begin{table}[htbp]
    \centering
    \caption{Work plan by year (Y1--Y3) and quarter (Q1--Q4).}
    \label{tab:gantt}
    \begin{tabular}{@{}llll@{}}
        \toprule
        Task & Year 1 & Year 2 & Year 3 \\
        \midrule
        Code design   & Q1--Q3 &        &        \\
        Decoder dev.  & Q2--Q4 & Q1--Q3 &        \\
        Simulation    & Q3--Q4 & Q1--Q4 & Q1--Q2 \\
        Hardware demo &        & Q3--Q4 & Q1--Q4 \\
        \bottomrule
    \end{tabular}
\end{table}

\chapter{Expected Outcomes}
\begin{enumerate}
    \item A published code family with a rigorous threshold
          proof~\cite{knuthTeXbook1986}.
    \item An open-source decoder with sub-microsecond inference on commodity
          GPUs~\cite{mckinneyPythonDataAnalysis2018}.
    \item A demonstration of logical qubit operation with error rate below
          $10^{-6}$ on a 72-qubit processor.
    \item At least three peer-reviewed publications and two conference talks.
\end{enumerate}

\chapter{Budget Justification}
The total requested budget is \$750\,000 over 36 months.
Table~\ref{tab:budget} provides the breakdown.

\begin{table}[htbp]
    \centering
    \caption{Budget summary by category.}
    \label{tab:budget}
    \begin{tabular}{@{}lr@{}}
        \toprule
        Category & Amount (\$) \\
        \midrule
        Personnel (2 PhD, 1 postdoc)  & 480\,000 \\
        Computing and cloud resources & 120\,000 \\
        Travel and conferences        &  45\,000 \\
        Equipment and supplies        &  65\,000 \\
        Indirect costs (15\%)         & 105\,000 \\
        \midrule
        \textbf{Total}               & \textbf{815\,000} \\
        \bottomrule
    \end{tabular}
\end{table}

\appendix
\chapter{Supporting Information}
Additional data, code listings, and supplementary derivations.

\printbibliography

\end{document}
